
\chapter{\texorpdfstring{$\zeta(s)$}{Zeta(s)} 与 \texorpdfstring{$\xi(s)$}{xi(s)} 的零点几何与分布}\label{ch:zero_distribution}

\section*{引言}\label{sec:intro_zero_distribution}

在延拓、函数方程与 $\xi$ 的 Hadamard 乘积之后，本章从几何上描述 $\zeta$ 与 $\xi$ 的零点。
目标：零点是实部与虚部零迹线的交点；临界带由已有定理限制；RH 在零迹线语言下有等价表述。

安排：先讨论 $\zeta$ 的实部、虚部曲面与零迹线（并与第~\ref{ch:riemann_zeta_continuation}~章模曲面对照）；
再给出 $N(T)$ 与临界带；最后讨论 $\xi$ 的对称曲面与临界线上 $\operatorname{Im}\xi\equiv0$。
图为示意，\textbf{不代替}零点区域定理与 RH 的证明。

\section{从 \texorpdfstring{$\xi(s)$}{xi(s)} 回到 \texorpdfstring{$\zeta(s)$}{Zeta(s)}}

上一章已构造 $\xi(s)$：它是 $\mathbb{C}$ 上的一阶整函数，并满足
\[
\xi(s)=\xi(1-s).
\]
$\xi$ 与 $\zeta$ 有完全相同的非平凡零点（见第~\ref{sec:sec-xi-zero-topology}~节）；
二者描述同一对象，规范化不同。
$\xi$ 为整函数且对称形式简单；$\zeta$ 则直接联系 Dirichlet 级数、Euler 乘积与素数分布。

第~\ref{ch:riemann_zeta_continuation}~章已给出全平面亚纯延拓、函数方程与分阶段模曲面示意。
本章在跨越临界带的绘图窗口上用实部、虚部曲面与零迹线表述零点，并讨论临界带内的区域限制与 $N(T)$。
图为示意，\textbf{不代替}证明。

\section{ Riemann \texorpdfstring{$\zeta$}{Zeta} 函数的实部曲面}
\label{sec:zeta_real_surface}

设 $s=\sigma+it$。将 $\zeta(s)$ 分解为
\[
\zeta(s)=\operatorname{Re}\zeta(s)+i\operatorname{Im}\zeta(s),
\]
则 $z=\operatorname{Re}\zeta(s)$ 是 $(\sigma,t)$ 平面上的实值曲面。

与第~\ref{ch:riemann_zeta_continuation}~章分阶段半平面图不同，本节取跨越临界带的窗口
$s\in[-4,4]\times[-32,32]$（覆盖 $0<\Re s<1$ 与 $\Re s=\tfrac12$）。
图~\ref{fig:zeta_real_cont} 给出解析延拓后的实部曲面，其上可见：
\begin{itemize}
    \item $s=1$ 邻域的一阶极点尖峰（$\zeta$ 的唯一极点；$\xi$ 无此结构）；
    \item $\Re(s)>1$ 上相对平缓，临界带附近振荡增强（Dirichlet 级数各项的相位叠加）；
    \item $\Re(s)<0$ 上振荡进一步加剧（与函数方程中 $\Gamma$、$\sin$ 因子一致）。
\end{itemize}
\begin{figure}[htbp]
\centering
\includegraphics[width=1.0\textwidth]{全书图形/11-Z实曲面.pdf}
\caption{解析延拓后的 $\zeta(s)$ 实部曲面。
参考面 $z=0$ 上：金黄色带为临界带 $0<\Re(s)<1$，橙色边线为带边界 $\Re(s)=0,1$；
蓝色虚线为临界线 $\Re(s)=\tfrac12$（函数方程对称轴，非 RH 证明）。
白色曲线为曲面与 $z=0$ 的交线（实部零迹线 $\Re\zeta=0$）。图为示意。}
\label{fig:zeta_real_cont}
\end{figure}

曲面与 $z=0$ 的交线 $\operatorname{Re}\zeta(s)=0$ 称为\textbf{实部零迹线}。
在临界带 $0<\sigma<1$ 内，它与虚部零迹线的交点对应非平凡零点；
左半实轴上的情形见下节与第~\ref{ch:riemann_zeta_continuation}~章图~\ref{fig:trivial_zeros}。
下节给出虚部曲面。

\section{ Riemann \texorpdfstring{$\zeta$}{Zeta} 函数的虚部曲面}
\label{sec:zeta_imag_surface}

同样，$z=\operatorname{Im}\zeta(s)$ 是 $(\sigma,t)$ 上的实值曲面。
在同一窗口 $s\in[-4,4]\times[-32,32]$ 上，图~\ref{fig:zeta_imag_cont} 给出虚部曲面。
由复共轭对称，虚部关于实轴 $t=0$ 呈奇对称，且在 $t=0$ 上恒为零。

由 Cauchy--Riemann 方程，实部与虚部局部相互联系。
交线 $\operatorname{Im}\zeta(s)=0$ 称为\textbf{虚部零迹线}；临界带内随 $t$ 增大振荡增强。

两曲面共用参考平面 $z=0$：白色零迹线的交点即为 $\zeta$ 的零点。
临界带内交点位置与第~\ref{ch:riemann_zeta_continuation}~章图~\ref{fig:zeta_modulus_cont_strip} 模曲面触底点一致；
左半实轴上交点落在 $s=-2,-4,-6$ 上（与图~\ref{fig:trivial_zeros} 及函数方程给出的平凡零点一致）。
零迹线交点与模曲面触底是同一事实的不同几何示意；图为示意，\textbf{不代替}证明。

图~\ref{fig:zeta_real_cont} 与图~\ref{fig:zeta_imag_cont} 是\textbf{分开展示}，分别读出 $\operatorname{Re}\zeta=0$ 与 $\operatorname{Im}\zeta=0$。
下一节将二者叠合，从「分镜」过渡到「合成」。

% 虚部曲面：取消浮动，固定于重叠图一节之前（避免漂到后文）
\begin{center}
\includegraphics[width=1.0\textwidth]{全书图形/11-Z虚曲面.pdf}
\captionof{figure}{解析延拓后的 $\zeta(s)$ 虚部曲面（共轭奇对称）。
参考面 $z=0$ 上：金黄色带为临界带 $0<\Re(s)<1$，橙色边线为 $\Re(s)=0,1$；
蓝色虚线为临界线 $\Re(s)=\tfrac12$（对称轴，非 RH 证明）。
白色曲线为虚部零迹线 $\Im\zeta=0$。图为示意。}
\label{fig:zeta_imag_cont}
\end{center}

\section{ Riemann \texorpdfstring{$\zeta$}{Zeta} 函数实部与虚部曲面的重叠图}

前两节已在同一窗口下分别给出实部、虚部曲面及其白色零迹线。
全纯函数的零点由两族零迹线的\textbf{公共}交点决定。叠合两曲面即可一次看清。

由定义 $\zeta(s)=\operatorname{Re}\zeta(s)+i\operatorname{Im}\zeta(s)$，
\[
\boxed{
\zeta(s)=0
\iff
\begin{cases}
\operatorname{Re}\zeta(s)=0,\\[2mm]
\operatorname{Im}\zeta(s)=0.
\end{cases}
}
\]
几何上：红色实部曲面与蓝色虚部曲面分别穿过 $z=0$，形成两组零迹线；公共交点即 $\zeta$ 的零点。

\begin{figure}[htbp]
\centering
\includegraphics[width=0.92\textwidth]{全书图形/Z曲面重叠.pdf}
\caption{$\zeta$ 实部与虚部曲面的重叠图（图~\ref{fig:zeta_real_cont} 与图~\ref{fig:zeta_imag_cont} 的\textbf{合成视角}）。
红：$\operatorname{Re}\zeta$；蓝：$\operatorname{Im}\zeta$；灰：参考平面 $z=0$。
公共交点满足 $\operatorname{Re}\zeta=\operatorname{Im}\zeta=0$，即 $\zeta$ 的非平凡零点。图为示意。}
\label{fig:zeta-real-imag-combined}
\end{figure}

图~\ref{fig:zeta-real-imag-combined} 将前两节两张曲面叠合：零点是 $z=0$ 上两族零迹线的交点。
「分镜」与「合成」读法一致，合成视角便于投影为二维零迹线网。

将临界带内零迹线投影到 $z=0$，得纵横交织的曲线网；每一交叉点对应一个非平凡零点。

\begin{figure}[htbp]
\centering
\includegraphics[width=0.90\textwidth]{全书图形/zeta零迹线.pdf}
\caption{$\zeta$ 零迹线在 $z=0$ 上的投影。红：$\operatorname{Re}\zeta=0$；蓝：$\operatorname{Im}\zeta=0$；绿：临界线 $\operatorname{Re}s=\tfrac12$。
纵轴 $\operatorname{Re}s$ 取负半轴朝上，与图~\ref{fig:zeta_real_cont}--\ref{fig:zeta-real-imag-combined} 的俯视投影一致（平凡零点在图上方）。
公共交点即非平凡零点。图为示意，\textbf{不代替}证明。}
\label{fig:zeta-zero-intersection}
\end{figure}

\subsection{零迹线相交的复分析约束——正交交汇}

临界带内红、蓝零迹线并非任意曲线，而受 \textbf{Cauchy--Riemann 方程} 制约。
设 $s=\sigma+it$，$\zeta(s)=u(\sigma,t)+iv(\sigma,t)$。
零迹线即隐曲线 $u=0$ 与 $v=0$。

在单零点 $s_0$（$\zeta'(s_0)\neq 0$）处，由
\[
u_{\sigma}=v_{t},\qquad u_{t}=-v_{\sigma}
\]
得 $v_{\sigma}=-u_{t}$、$v_{t}=u_{\sigma}$，从而
\[
\nabla u\cdot\nabla v
=u_{\sigma}v_{\sigma}+u_{t}v_{t}
=u_{\sigma}(-u_{t})+u_{t}u_{\sigma}=0.
\]
故 $\nabla u\perp\nabla v$。零迹线的切向分别垂直于 $\nabla u$、$\nabla v$，
因而两族切向亦正交：在 $z=0$ 上，单零点处两族零迹线以直角相交。
图~\ref{fig:zeta-zero-intersection} 与此一致。

比较图~\ref{fig:zeta-real-imag-combined} 与图~\ref{fig:zeta-zero-intersection}：
前者为全局曲面，后者为零迹线投影。
下一节讨论临界带内的区域限制与计数。

\section{非平凡零点的分布与临界带}\label{sec:zero-distribution-strip}

前几节已建立几何语言：零点是实部与虚部零迹线的公共交点。
问题转为区域限制：非平凡零点落在复平面的何处？

第~\ref{ch:riemann_zeta_continuation}~章已由 Euler 乘积与函数方程给出初步结论；
本节在零迹线语言下复述并收紧（含 $\operatorname{Re}s=1$ 上无零点），并引入 $N(T)$。

当 $\operatorname{Re}(s)>1$ 时，$\zeta$ 有绝对收敛的 Euler 乘积
\[
\zeta(s)=\prod_{p}\frac{1}{1-p^{-s}}.
\]
每一因子非零，故乘积非零：
\[
\boxed{
\operatorname{Re}(s)>1
\Longrightarrow
\zeta(s)\neq0.
}
\]

另一方面，函数方程
\[
\zeta(s)
=
2^s\pi^{s-1}\sin\frac{\pi s}{2}\,\Gamma(1-s)\,\zeta(1-s)
\]
在 $\operatorname{Re}(s)<0$ 时给出 $\operatorname{Re}(1-s)>1$，从而 $\zeta(1-s)\neq 0$。
因子 $2^s\pi^{s-1}$ 恒非零；$\Gamma(1-s)$ 无零点（$1/\Gamma$ 为整函数）。
故左半平面零点只能来自 $\sin(\pi s/2)=0$，即 $s=-2,-4,-6,\ldots$：
\[
\boxed{
\operatorname{Re}(s)<0
\text{ 的零点仅有 }
-2,-4,-6,\ldots
}
\]
这些称为\textbf{平凡零点}。

于是右半平面与左半平面的零点结构已定；可能出现非平凡零点的区域只剩
\[
0<\operatorname{Re}(s)<1,
\]
称为\textbf{临界带}。

Hadamard 与 de la Vallée Poussin（1896）独立证明
\[
\boxed{
\zeta(s)\neq0,\qquad\operatorname{Re}(s)=1.
}
\]
故零点不能落在临界带右边界。结合 $\zeta(s)=\overline{\zeta(\bar s)}$ 与函数方程，亦得 $\operatorname{Re}(s)=0$ 上无非平凡零点。
从而非平凡零点严格位于开带 $0<\operatorname{Re}(s)<1$。

\subsection{零点计数函数 \texorpdfstring{$N(T)$}{N(T)} 与 Riemann--von Mangoldt 公式}
\label{sec:NT-RVM}

为定量描述临界带内零点的垂直分布，定义\textbf{零点计数函数}
\[
N(T)
=
\#\bigl\{\rho=\beta+i\gamma:\;
0<\beta<1,\;
0<\gamma\le T
\bigr\}
\]
（计重数；仅上半平面）。$N(T)$ 是 $T$ 的\textbf{右连续阶梯函数}：
在非零点纵坐标处局部常数，每经过一个虚部 $\gamma$（计重数）跳跃 $1$。
图~\ref{fig:NT-staircase} 取 $0\le T\le 80$，与书中相应三维示意动画的截断高度一致；
$N(T)$ 在每一枚 $\gamma_n\le 80$ 处跳一阶（共 $21$ 级）。

\begin{figure}[htbp]
\centering
\begin{tikzpicture}[x=0.12cm,y=0.28cm,
  font=\small,
  >=Stealth]
  % ---------- 网格（先画）----------
  \draw[xstep=5, ystep=1, gray!15, very thin] (0,0) grid (80,21);
  % 水平线强调计数轴 N=1,2,...,21
  \foreach \y in {1,...,21}
    \draw[DarkGreen!25, thin] (0,\y) -- (80,\y);
  \foreach \x in {0,10,20,30,40,50,60,70,80}
    \draw[gray!40, thin] (\x,0) -- (\x,21);
  \foreach \y in {0,5,10,15,20}
    \draw[gray!45, thin] (0,\y) -- (80,\y);

  % 坐标轴
  \draw[->] (0,0) -- (88,0) node[right] {$T$};
  \draw[->,DarkGreen!70!black] (0,0) -- (0,24)
    node[above,DarkGreen!70!black] {$N(T)$~~\textnormal{\footnotesize（计数轴）}};
  \foreach \x in {0,10,20,30,40,50,60,70,80}
    \draw (\x,0.25) -- (\x,-0.25) node[below,font=\footnotesize] {$\x$};
  % 纵轴计数点 N=0,1,...,21
  \foreach \y in {0,1,...,21} {
    \draw[DarkGreen!70!black] (0.35,\y) -- (-0.35,\y);
  }
  \foreach \y in {0,5,10,15,20,21}
    \node[left,font=\footnotesize,DarkGreen!70!black] at (-0.45,\y) {$\y$};
  \foreach \y in {1,2,3,4,6,7,8,9,11,12,13,14,16,17,18,19} {
    \node[left,font=\scriptsize,DarkGreen!55!black] at (-0.35,\y) {$\y$};
  }

  \def\zeroList{%
      14.13/1,21.02/2,25.01/3,30.42/4,32.94/5,37.59/6,40.92/7,43.33/8,%
      48.01/9,49.77/10,52.97/11,56.45/12,59.35/13,60.83/14,65.11/15,%
      67.08/16,69.55/17,72.07/18,75.70/19,77.14/20,79.34/21}

  \foreach \g/\n in \zeroList {
    \draw[DarkGreen!60, densely dotted, thin] (0,\n) -- (\g,\n);
    \draw[gray!55, densely dotted, thin] (\g,0) -- (\g,\n);
    \fill[DarkGreen] (0,\n) circle (1.6pt);
    \draw[DarkGreen!80!black, thin] (0,\n) circle (2.2pt);
  }

  % 右连续阶梯
  \draw[thick,MainBlue,line cap=round]
    (0,0) -- (14.13,0)
    -- (14.13,1) -- (21.02,1)
    -- (21.02,2) -- (25.01,2)
    -- (25.01,3) -- (30.42,3)
    -- (30.42,4) -- (32.94,4)
    -- (32.94,5) -- (37.59,5)
    -- (37.59,6) -- (40.92,6)
    -- (40.92,7) -- (43.33,7)
    -- (43.33,8) -- (48.01,8)
    -- (48.01,9) -- (49.77,9)
    -- (49.77,10) -- (52.97,10)
    -- (52.97,11) -- (56.45,11)
    -- (56.45,12) -- (59.35,12)
    -- (59.35,13) -- (60.83,13)
    -- (60.83,14) -- (65.11,14)
    -- (65.11,15) -- (67.08,15)
    -- (67.08,16) -- (69.55,16)
    -- (69.55,17) -- (72.07,17)
    -- (72.07,18) -- (75.70,18)
    -- (75.70,19) -- (77.14,19)
    -- (77.14,20) -- (79.34,20)
    -- (79.34,21) -- (80,21);

  \foreach \g/\n in \zeroList {
    \fill[MainBlue] (\g,\n) circle (1.35pt);
  }

  \draw[thick,dashed,DarkRed,domain=18:80,samples=100,smooth]
    plot (\x, { max(0, (\x/(2*pi))*ln(\x/(2*pi*2.71828)) ) });
  \node[DarkRed,anchor=west,font=\footnotesize] at (62,12.5)
    {$\frac{T}{2\pi}\log\frac{T}{2\pi e}$};

  \node[MainBlue,anchor=west,font=\footnotesize] at (2,23.0)
    {$N(T)$ 阶梯};
  \node[DarkGreen,anchor=west,font=\footnotesize] at (22,23.0)
    {\textbullet~\textbf{计数点} $N=1,2,\ldots$（纵轴）};
\end{tikzpicture}
\caption{零点计数函数 $N(T)$ 的阶梯示意（$0\le T\le 80$）。
\textbf{纵轴}为计数值 $N(T)$：每个整数高度 $N=1,2,\ldots,21$ 为一级\textbf{计数点}
（绿色圆点及水平网格线），阶梯在 $\gamma_n$ 处跳到下一计数点。
横轴为高度 $T$；浅色网格便于读 $T$ 与 $N$。
本图取 $\gamma_n\le 80$ 的 $21$ 个纵坐标，与相应三维示意动画中
高度截断 $T\le 80$ 时零点计数器的读数一致。
红色虚线为 Riemann--von Mangoldt 主项示意。}
\label{fig:NT-staircase}
\end{figure}

经典的 Riemann--von Mangoldt 公式断言
\[
\boxed{
N(T)
=
\frac{T}{2\pi}\log\frac{T}{2\pi e}
+
O(\log T)
\qquad(T\to\infty).
}
\]
主项来自 $\xi$ 在临界线上的辐角变化（Stirling 控制 $\Gamma$ 因子）；$O(\log T)$ 为经典误差界，可改进。
相邻零点平均间距
\[
\Delta\gamma\sim\frac{2\pi}{\log(T/2\pi)}.
\]
图~\ref{fig:NT-staircase} 中台阶随 $T$ 变密，与上式一致。
本章后文几何图像与第~\ref{ch:fourier_vs_riemann}~章频谱比较均依赖这一计数框架；
完整证明属整函数增长理论，陈述见 Titchmarsh 等标准文献。

\begin{historicalnote}
\textbf{历史说明\quad Hadamard 与 de la Vallée Poussin。}
1896 年，二人独立证明 $\zeta(s)\neq0$（$\operatorname{Re}(s)=1$），
排除零点落在临界带右边界，并几乎同时完成素数定理 $\pi(x)\sim x/\log x$。
证明综合整函数理论、Hadamard 乘积与对数导数估计。
\textbf{证明提要见第~\ref{ch:prime_number_theorem}~章}；
$3$--$4$--$1$ 不等式与零点自由区见附录~\ref{app:pnt_estimates}。
亦可参阅 Titchmarsh \textit{The Theory of the Riemann Zeta-Function}
与 Edwards \textit{Riemann's Zeta Function}。
\end{historicalnote}

综合得 $\zeta$ 零点的区域限制：
\[
\boxed{
\begin{array}{ll}
\operatorname{Re}(s)>1, & \zeta(s)\neq0, \\[2mm]
\operatorname{Re}(s)=1, & \zeta(s)\neq0, \\[2mm]
\operatorname{Re}(s)<0, & \text{仅有平凡零点 }-2,-4,-6,\ldots, \\[2mm]
0<\operatorname{Re}(s)<1, & \text{全部非平凡零点所在区域。}
\end{array}
}
\]
前节零迹线重叠图与此一致：公共交点均落在临界带内（图为示意，\textbf{不代替}证明）。

Riemann（1859）进一步问：非平凡零点是否全部位于 $\operatorname{Re}(s)=1/2$？
此即下一节的黎曼猜想（RH）。

\section{临界直线与黎曼猜想}\label{sec:critical_line_rh}

非平凡零点已限制在开临界带 $0<\operatorname{Re}(s)<1$，但带宽为 $1$ 的区域内精细分布仍开放。

1859 年，Riemann 在《论小于给定量的素数个数》中提出：
\begin{center}
\textbf{$\zeta$ 的全部非平凡零点都位于直线 $\operatorname{Re}(s)=\frac12$ 上。}
\end{center}
即\textbf{黎曼猜想（RH）}：
\[
\boxed{
\operatorname{Re}(\rho)=\frac12,\quad \forall \rho\text{ 满足 }\zeta(\rho)=0\text{ 且 }0<\operatorname{Re}(\rho)<1.
}
\]

这是关于全部非平凡零点的\textbf{全称命题}：
一枚 $\operatorname{Re}(\rho)\neq 1/2$ 的零点即可否证；
任意有限高度上的确认只给出「虚部 $\le T$ 的零点均在临界线上」的\textbf{有限全称}结论，
\textbf{不能}代替无限多个零点的证明。
正因一反例即可否证，在可达到高度上作可严格化的大规模计算具有理论必要性：
搜寻或排除到某高度的反例，并提供该高度内可引用的有限结论
（见第~\ref{ch:riemann_numerical_approximation}、\ref{ch:contemporary_research}~章）。

由 $\xi(s)=\xi(1-s)$，若 $\rho$ 为非平凡零点，则 $1-\rho$、$\bar\rho$、$1-\bar\rho$ 亦然；
故非平凡零点关于 $\operatorname{Re}(s)=1/2$ 与实轴镜像对称。
RH 断言该对称结构退化为全部零点落在同一条竖直直线上。

几何上：实部零迹线 $\operatorname{Re}\zeta=0$ 与虚部零迹线 $\operatorname{Im}\zeta=0$ 在临界带内交织，
每一非平凡零点是其交点；RH 断言全部交点集中在 $\operatorname{Re}(s)=1/2$。

RH 尚未证明。数值上，虚部不超过约 $3\times 10^{12}$ 的非平凡零点（约 $10^{13}$ 量级）
均已验证落在临界线上（如 Platt \& Trudgian, 2021），未见反例——
此为有限高度内的确认，不是 RH 的证明。

RH 是 $\zeta$ 理论的中心问题之一，与素数分布及诸多解析数论结果相关；
2000 年 Clay 数学研究所将其列入七个 \emph{Millennium Prize Problems} 之一。
若 RH 成立，两族零迹线的全部公共交点均在 $\operatorname{Re}(s)=1/2$ 上，
显式公式中零点项相位将受精确控制，从而改进素数计数的误差估计。

\section{两种几何语言的汇合：从 \texorpdfstring{$\zeta$}{Zeta} 回到 \texorpdfstring{$\xi$}{xi}}\label{sec:two-languages}

从 $\zeta$ 视角已得：非平凡零点是临界带内两族零迹线的交点；RH 断言它们全部位于 $\operatorname{Re}(s)=1/2$。

$\zeta$ 在 $s=1$ 有极点，函数方程含 $\Gamma$ 与三角函数，对称 $s\mapsto 1-s$ 形式不够直接。
$\xi$ 消去极点与附加因子后为整函数且 $\xi(s)=\xi(1-s)$，非平凡零点不变。

二者描述同一组非平凡零点，侧重点不同：
\begin{itemize}
  \item \textbf{$\zeta$}：联系 Dirichlet 级数、Euler 乘积与素数分布；素数定理与显式公式的自然出发点。
  \item \textbf{$\xi$}：Hadamard 乘积将函数编码为零点集合；曲面光滑对称，临界线作为对称轴更易辨认。
\end{itemize}

\noindent\textbf{叙述安排：}前半用 $\zeta$ 确立区域限制与 RH 的经典陈述；后半用 $\xi$ 的零迹线给出 RH 的几何等价表述。

\section{\texorpdfstring{$\xi$}{xi} 曲面的几何结构}

上一章已用 Hadamard 乘积表明：在适当正规化下，
$\zeta$ 的非平凡零点（亦即 $\xi$ 的全部零点）唯一确定整函数 $\xi$。
本节考察 $\xi$ 的几何形态：实部、虚部曲面及其与 $z=0$ 的交线（零迹线）；
实部曲面与虚部曲面的零迹线交点给出同一零点集。

设 $s=\sigma+it$。$\xi(s)=\operatorname{Re}\xi(s)+i\operatorname{Im}\xi(s)$ 对应两张实值曲面
$z=\operatorname{Re}\xi(s)$ 与 $z=\operatorname{Im}\xi(s)$。
\begin{figure}[htbp]
\centering
\includegraphics[width=0.90\textwidth]{全书图形/xi实虚曲面重叠.pdf}
\caption{$\xi$ 实部与虚部曲面的重叠图。
红：$\operatorname{Re}\xi$；蓝：$\operatorname{Im}\xi$；灰：$z=0$。
关于 $\operatorname{Re}(s)=1/2$ 镜像对称；交线即两族零迹线，公共交点为 $\xi$ 的全部零点
（即 $\zeta$ 的全部非平凡零点）。图为示意。}
\label{fig:xi-real-imag-combined}
\end{figure}

图~\ref{fig:xi-real-imag-combined} 显示：$\xi$ 已消去 $\zeta$ 的极点与 $\Gamma$ 因子奇性，曲面光滑且关于
$\operatorname{Re}(s)=1/2$ 镜像对称；$\operatorname{Re}\xi=0$ 与 $\operatorname{Im}\xi=0$ 的交点即全部非平凡零点。

下一节讨论零迹线拓扑，并将 RH 表述为交汇规律的几何命题。

\section{零迹线拓扑与黎曼猜想的几何视角}\label{sec:sec-xi-zero-topology}

由
\[
\xi(s)
=
\frac12 s(s-1)\pi^{-s/2}\Gamma\!\left(\frac{s}{2}\right)\zeta(s),
\]
除 $\zeta$ 外其余因子不产生非平凡零点，故
\[
\boxed{
\xi(s)\text{ 与 }\zeta(s)\text{ 具有完全相同的非平凡零点。}
}
\]
研究 $\zeta$ 的非平凡零点等价于研究整函数 $\xi$ 的零点。
$\xi$ 无极点，满足 $\xi(s)=\xi(1-s)$，关于 $\operatorname{Re}(s)=1/2$ 对称更清晰。

\begin{figure}[htbp]
\centering
\includegraphics[width=0.90\textwidth]{全书图形/xi零迹线.pdf}
\caption{$\xi$ 零迹线在 $z=0$ 上的投影。
红：$\operatorname{Re}\xi=0$；蓝：$\operatorname{Im}\xi=0$；绿：临界线 $\operatorname{Re}s=\tfrac12$
（亦为 $\operatorname{Im}\xi=0$ 的一条连通分支，见下文推导）。
纵轴 $\operatorname{Re}s$ 取负半轴朝上，与图~\ref{fig:zeta_real_cont}--\ref{fig:zeta-zero-intersection} 一致。
公共交点即非平凡零点。图为示意，\textbf{不代替}证明。}
\label{fig:xi-real-imag-intersection}
\end{figure}

图~\ref{fig:xi-real-imag-intersection} 将三维结构投影为二维零迹线网。
临界线上 $\operatorname{Im}\xi\equiv 0$ 可分两步得到。

\textbf{第一步（函数方程）。}由 $\xi(s)=\xi(1-s)$，令 $s=\tfrac12+it$，则
\[
\xi\!\left(\frac12+it\right)=\xi\!\left(\frac12-it\right).
\]

\textbf{第二步（Schwarz 反射）。}$\xi$ 在实轴上取实值，故 $\overline{\xi(\bar s)}=\xi(s)$。
令 $s=\tfrac12+it$，得
\[
\overline{\xi\!\left(\frac12+it\right)}=\xi\!\left(\frac12-it\right).
\]
综合两步：$\xi(1/2+it)=\overline{\xi(1/2+it)}$，即
\[
\boxed{
\operatorname{Im}\xi\!\left(\frac12+it\right)=0.
}
\]
于是整条临界线
\[
\boxed{\operatorname{Re}(s)=\frac12}
\]
构成 $\operatorname{Im}\xi=0$ 的一条连通分支。

图~\ref{fig:xi-real-imag-intersection} 中，该分支贯穿临界带；
$\operatorname{Re}\xi=0$ 为另一族零迹线；交点即全部非平凡零点。
从而 RH 可表述为：临界带内两族零迹线的全部交点都落在 $\operatorname{Re}(s)=1/2$ 上
（等价重述，不是新证明）。

数值与图形观察还提示：除临界线外，其余 $\operatorname{Im}\xi=0$ 分支在临界带内
似乎不与 $\operatorname{Re}\xi=0$ 相交。若能严格证明这一点，则全部非平凡零点只能位于
$\operatorname{Re}(s)=1/2$，RH 立即成立。
此仍属零迹线拓扑层面的几何观察，连通性与整体拓扑尚需严格理论；
它是 RH 的一种几何等价表述，不是新证明。

\begin{center}
\fbox{\parbox{0.88\textwidth}{\centering
RH 可理解为 $\operatorname{Re}\xi=0$ 与 $\operatorname{Im}\xi=0$
两族零迹线拓扑结构的几何命题。}}
\end{center}

\begin{center}
\fbox{\parbox{0.92\textwidth}{
\centering
\vspace{0.2em}
\textbf{第十一章：重要定理与核心公式汇总}\\[6pt]
\renewcommand{\arraystretch}{1.6}
\begin{tabular}{lp{9.5cm}}
\hline
\textbf{名称} & \textbf{数学内容 / 公式表达} \\
\hline
$\zeta$ 实部 / 虚部曲面 & 极点尖峰、振荡；零迹线 $\Re\zeta=0$、$\Im\zeta=0$ \\
$\xi$ 实虚重叠与零迹线 & 关于 $\Re s=1/2$ 对称；交点即非平凡零点 \\
零迹线正交交汇 & 单零点处 $\nabla u\perp\nabla v$（Cauchy--Riemann） \\
Riemann--von Mangoldt 公式 & $N(T)=\dfrac{T}{2\pi}\log\dfrac{T}{2\pi e}+O(\log T)$ \\
黎曼猜想几何等价 & 临界带内两族零迹线交点均在 $\Re s=\tfrac12$ \\
\hline
\end{tabular}
\vspace{0.2em}
}}
\end{center}

\section*{本章小结与后续展望}\label{sec:zero_distribution_summary}
\addcontentsline{toc}{section}{本章小结与后续展望}

本章用曲面与零迹线描述 $\zeta$ 与 $\xi$ 的零点：交点对应零点；临界带由 Euler 乘积、函数方程与 $\operatorname{Re}s=1$ 无零点确定；$N(T)$ 给出垂直计数。对 $\xi$，对称性使 $\operatorname{Re}s=1/2$ 成为 $\operatorname{Im}\xi=0$ 的一条分支；RH 在几何上等价于：临界带内实部与虚部零迹线的全部交点都落在该轴上（这是等价重述，不是新证明）。

下一章（第~\ref{ch:prime_number_theorem}~章）给出素数定理的证明提要：
在 $\operatorname{Re}s=1$ 上无零点的条件下建立 $\psi(x)\sim x$，从而 $\pi\sim\Li$。
再下一章（第~\ref{ch:riemann_explicit_formula}~章）用围道积分建立 $\psi_0$ 的显式公式，
把非平凡零点与素数计数函数直接联系起来。